\documentclass[a4paper,11pt]{ltjsarticle}

\usepackage[margin=23mm]{geometry}
\usepackage{amsmath,amssymb,amsthm,mathtools}
\usepackage{booktabs}
\usepackage{xcolor}
\usepackage[hidelinks]{hyperref}
\usepackage{enumitem}

\definecolor{ink}{HTML}{1F2933}
\definecolor{accent}{HTML}{6B4F9A}
\definecolor{soft}{HTML}{F4F1FA}

\hypersetup{
  colorlinks=true,
  linkcolor=accent,
  urlcolor=accent,
  citecolor=accent
}

\setlist[itemize]{leftmargin=1.5em,itemsep=0.22em,topsep=0.25em}
\setlength{\parindent}{1em}
\setlength{\parskip}{0.18em}

\theoremstyle{definition}
\newtheorem{definition}{定義}[section]
\newtheorem{example}{例}[section]

\theoremstyle{plain}
\newtheorem{proposition}{命題}[section]
\newtheorem{theorem}{定理}[section]

\theoremstyle{remark}
\newtheorem*{remark}{注意}

\newcommand{\softmax}{\operatorname{softmax}}
\newcommand{\lse}{\operatorname{lse}}
\newcommand{\Att}{\operatorname{Att}}

\title{\vspace{-1.1cm}
{\Large\color{accent}Memory and Temperature}\\[0.35em]
{\huge Hopfield Networks と Temperature Sampling}\\[0.35em]
{\large Hanoi 実験を読むための短い教科書}}
\author{}
\date{}

\begin{document}
\color{ink}
\maketitle
\vspace{-0.9cm}

\begin{center}
\begin{minipage}{0.88\linewidth}
\small
本ノートは，Ramsauer et al. の
\textit{Hopfield Networks is All You Need}
と，Wu--Mirhoseini--Tambe の
\textit{On the Role of Temperature Sampling in Test-Time Scaling}
を，発表で使う範囲に絞って整理する。
前者は Transformer attention を連想記憶の検索として見る理論を与え，
後者は test-time scaling をサンプル数だけでなく温度方向にも広げる。
Hanoi 実験では，低温での反復的失敗と高温での回復を説明する足場になる。
\end{minipage}
\end{center}

\section{二つの論文をつなぐ見方}

LLM の推論を，次の二段階で見る。

\begin{itemize}
\item 内部では，attention が文脈中の情報・中間表現・プロトタイプを検索する。
\item 出力では，logit 分布から次トークンをサンプリングし，推論軌道を進める。
\end{itemize}

Ramsauer et al. は第一の段階に理論的な言葉を与える。
Transformer attention は，modern Hopfield network の更新則と対応する。
つまり attention は単なる重み付き平均ではなく，記憶パターンへの連想検索として読める。

Wu et al. は第二の段階に焦点を当てる。
同じモデル，同じ問題でも，温度を変えると生成される reasoning trace の集合が変わる。
そのため，単一温度でサンプル数 \(K\) を増やすだけでは到達できない解が，
別の温度で現れることがある。

\begin{center}
\begin{tabular}{lll}
\toprule
論文 & 主対象 & 発表での役割 \\
\midrule
Hopfield Networks is All You Need & attention と連想記憶 & 失敗を attractor として語る \\
Temperature Sampling in TTS & 温度と推論軌道 & 高温で別解法に届く可能性を語る \\
\bottomrule
\end{tabular}
\end{center}

\section{Modern Hopfield Network}

古典的 Hopfield network は，記憶したパターンをエネルギーの局所最小点として持つ
連想記憶モデルである。Ramsauer et al. の modern Hopfield network は，
これを連続状態に拡張し，指数的な記憶容量と attention との対応を示した。

\begin{definition}[連想記憶としての検索]
記憶パターンを列ベクトルとして並べた行列を
\[
  X=(\xi_1,\ldots,\xi_N)\in\mathbb{R}^{d\times N}
\]
とする。状態 \(x\in\mathbb{R}^d\) が与えられたとき，modern Hopfield update は
\[
  x^{\mathrm{new}}
  =
  X\,\softmax(\beta X^\top x)
\]
で与えられる。ここで \(\beta>0\) は inverse temperature に相当し，
内積 \(X^\top x\) は状態 \(x\) と各記憶パターンの類似度を表す。
\end{definition}

\(\softmax(\beta X^\top x)\) は，どの記憶パターンをどれだけ読むかを決める重みである。
\(\beta\) が大きいほど最も近い記憶に鋭く集中し，小さいほど複数の記憶を平均する。

\begin{definition}[エネルギーの形]
論文では，状態 \(x\) に対して log-sum-exp を含むエネルギーを考える。
概念的には
\[
  E(x)
  =
  -\frac{1}{\beta}\log\sum_{i=1}^{N}
  \exp(\beta \xi_i^\top x)
  +\frac{1}{2}\|x\|^2
  +\text{定数}
\]
という形である。更新則はこのエネルギーを下げる方向の写像として理解できる。
\end{definition}

\begin{theorem}[Ramsauer et al. の主結果]
Modern Hopfield network は，連想空間の次元に対して指数的に多くのパターンを保存でき，
一回の更新で記憶を検索でき，retrieval error は指数的に小さくなる。
また，その更新則は Transformer の attention と等価に解釈できる。
\end{theorem}

論文は energy minima を三種類に分ける。

\begin{center}
\begin{tabular}{ll}
\toprule
固定点の型 & 意味 \\
\midrule
global averaging & 多くのパターンの平均へ落ちる \\
metastable state & 一部のパターン集合の平均へ落ちる \\
single-pattern state & ほぼ一つの記憶を検索する \\
\bottomrule
\end{tabular}
\end{center}

\section{Attention との対応}

Transformer attention は，query \(Q\)，key \(K\)，value \(V\) に対して
\[
  \Att(Q,K,V)
  =
  \softmax\!\left(\frac{QK^\top}{\sqrt{d}}\right)V
\]
と書ける。これは，query と key の類似度を softmax で重みに変え，
value の重み付き平均を返す操作である。

Modern Hopfield update
\[
  x^{\mathrm{new}}
  =
  X\,\softmax(\beta X^\top x)
\]
と比べると，query が現在状態，key が検索対象，value が読み出される内容に対応する。
したがって attention head は，文脈中のトークンや内部表現から，
必要な記憶を検索する Hopfield 的な層として読める。

\begin{remark}
この論文が直接主張しているのは，
attention の数式が modern Hopfield update と対応するということである。
「LLM が本当に人間の記憶のように考えている」という主張ではない。
発表では，attention を記憶検索として解釈できる，という範囲に留めるのがよい。
\end{remark}

Hanoi へつなぐなら，正しい再帰解法を一つの安定した検索先，
反復的な誤手順を別の安定した検索先として見る。
低温で同じ誤手順に落ち続けるなら，それは出力分布上の高確率な局所軌道であり，
内部表現としては spurious attractor のように振る舞っている，と説明できる。

\section{Temperature Sampling と Test-Time Scaling}

次に，出力トークンの選び方を見る。
言語モデルは各時刻で logit \(z_i\) を出し，温度 \(T>0\) を使って
\[
  p_T(i)
  =
  \frac{\exp(z_i/T)}
       {\sum_j \exp(z_j/T)}
\]
から次トークンを選ぶ。

\begin{itemize}
\item \(T\to0\): 最大 logit のトークンに集中する。決定的で反復しやすい。
\item \(T=1\): 元の確率分布に従う。
\item \(T>1\): 分布が平らになり，低確率候補も選ばれやすくなる。
\end{itemize}

Test-time scaling (TTS) では，同じ問題に対して複数の reasoning trace を生成し，
答えを選択・投票・検証する。単一温度なら
\[
  \mathcal{R}_{T,K}
  =
  \{r_1,\ldots,r_K\},\qquad
  r_i\sim p_T(r\mid x)
\]
のように，温度 \(T\) を固定して \(K\) 本の推論を生成する。

\begin{theorem}[Wu et al. の主張]
単一温度でサンプル数 \(K\) を増やすだけでは，性能向上は無限には続かない。
一方，異なる温度は異なる問題集合を解くため，温度方向に scaling すると，
単一温度 TTS より広い reasoning boundary を探索できる。
\end{theorem}

論文は Qwen3 系列の複数サイズと，AIME 2024/2025，MATH500，
LiveCodeBench，Hi-ToM などの reasoning benchmark で評価している。
平均すると，温度方向の scaling は単一温度 TTS に対して追加の性能改善を与えた。
さらに，multi-temperature voting により，複数温度を使う計算コストを抑える方法も提案している。

\begin{definition}[multi-temperature voting の考え方]
温度集合 \(\mathcal{T}=\{T_1,\ldots,T_m\}\) を用意し，各温度で少数の推論を生成する。
\[
  \mathcal{R}
  =
  \bigcup_{T\in\mathcal{T}}
  \{r_{T,1},\ldots,r_{T,k_T}\}.
\]
得られた答えを集約し，多数決や検証器で最終答えを選ぶ。
これは，サンプル数 \(K\) だけでなく温度軸にも計算資源を配る方法である。
\end{definition}

\section{Hanoi 実験での読み方}

Hanoi では，低温で正答率が低く，高温で一部の \(N\) だけ正答率が上がることがある。
この現象は，次のように読むと二つの論文を自然につなげられる。

\begin{itemize}
\item 低温では，モデルは高確率な定型手順を選び続ける。
\item その定型手順が正しい再帰解法なら成功するが，誤った move loop なら失敗が固定化する。
\item 高温では，次トークン選択が平らになり，別の reasoning trace に入る確率が上がる。
\item その別軌道が正しい解法アトラクタに近ければ，正答率が回復する。
\item ただし高すぎる温度では，軌道そのものが壊れ，合法手や一貫性が失われる。
\end{itemize}

この説明で重要なのは，高温を単なるノイズとして扱わないことである。
Wu et al. の結果に従えば，温度は探索する推論経路の集合を変える操作である。
また Ramsauer et al. の見方に従えば，attention は記憶・プロトタイプ・中間表現への検索として読める。

\begin{center}
\begin{tabular}{lll}
\toprule
観測 & Hopfield 的解釈 & temperature 的解釈 \\
\midrule
同じ誤手順の反復 & spurious attractor & 低温で高確率軌道に固定 \\
高温で正答率上昇 & 別 basin への遷移 & 異なる reasoning trace の探索 \\
高温で崩壊 & 検索先が不安定化 & 分布が平らすぎて制御不能 \\
\bottomrule
\end{tabular}
\end{center}

\section*{発表での一言}

Transformer attention は，modern Hopfield network と対応するため，
記憶検索や attractor の言葉で解釈できる。
一方，temperature sampling は，その内部表現から出力軌道を選ぶ段階で，
どの reasoning trace を探索するかを変える。
したがって Hanoi の高温回復は，
「ノイズで偶然よくなった」というより，
「低温で固定された誤った軌道から外れ，別の解法軌道に入った」
可能性として話せる。

\section*{参考文献}

\begin{enumerate}[leftmargin=1.7em]
\item Hubert Ramsauer et al.,
``Hopfield Networks is All You Need,'' ICLR 2021. arXiv:2008.02217.\\
\url{https://arxiv.org/abs/2008.02217}
\item Yuheng Wu, Azalia Mirhoseini, Thierry Tambe,
``On the Role of Temperature Sampling in Test-Time Scaling,'' arXiv:2510.02611, 2025.\\
\url{https://arxiv.org/abs/2510.02611}
\end{enumerate}

\end{document}
